\documentclass[12pt, a4paper]{article}

% ─── Paquetes ────────────────────────────────────────────────────────────────
\usepackage[utf8]{inputenc}
\usepackage[T1]{fontenc}
\usepackage[spanish]{babel}
\usepackage{geometry}
\usepackage{titlesec}
\usepackage{fancyhdr}
\usepackage{booktabs}
\usepackage{array}
\usepackage{tabularx}
\usepackage{xcolor}
\usepackage{mdframed}
\usepackage{enumitem}
\usepackage{hyperref}
\usepackage{parskip}
\usepackage{lmodern}
\usepackage{microtype}
\usepackage{float}
\usepackage{amssymb}

% ─── Geometría ───────────────────────────────────────────────────────────────
\geometry{top=2.5cm, bottom=2.5cm, left=3cm, right=2.5cm, headheight=15pt}

% ─── Colores ─────────────────────────────────────────────────────────────────
\definecolor{primary}{HTML}{1A56DB}
\definecolor{lightblue}{HTML}{EBF5FF}
\definecolor{darkgray}{HTML}{374151}
\definecolor{accentgreen}{HTML}{046C4E}
\definecolor{lightgreen}{HTML}{F3FAF7}
\definecolor{accentorange}{HTML}{B45309}
\definecolor{lightorange}{HTML}{FFFBEB}
\definecolor{bordercolor}{HTML}{D1D5DB}
\definecolor{sectionbg}{HTML}{F9FAFB}
\definecolor{accentred}{HTML}{9B1C1C}
\definecolor{lightred}{HTML}{FDF2F2}
\definecolor{scoreok}{HTML}{D1FAE5}
\definecolor{scorewarn}{HTML}{FEF3C7}

% ─── Encabezado y pie de página ───────────────────────────────────────────────
\pagestyle{fancy}
\fancyhf{}
\renewcommand{\headrulewidth}{0.5pt}
\renewcommand{\headrule}{\hbox to\headwidth{\color{primary}\leaders\hrule height \headrulewidth\hfill}}
\fancyhead[L]{\small\color{darkgray}\textbf{TaskFlow} --- Informe de Accesibilidad}
\fancyhead[R]{\small\color{darkgray}v1.0 · 24 de junio de 2026}
\fancyfoot[C]{\small\color{darkgray}\thepage}

% ─── Títulos de sección ───────────────────────────────────────────────────────
\titleformat{\section}[block]
  {\large\bfseries\color{primary}}
  {}
  {0pt}
  {\colorbox{lightblue}{\parbox{\dimexpr\linewidth-2\fboxsep}{\strut #1\strut}}}

\titleformat{\subsection}[block]
  {\normalsize\bfseries\color{darkgray}}
  {}
  {0pt}
  {#1}

\titleformat{\subsubsection}[runin]
  {\small\bfseries\color{darkgray}}
  {}
  {0pt}
  {#1\quad}

% ─── Entornos personalizados ──────────────────────────────────────────────────
\newmdenv[
  backgroundcolor=lightblue,
  linecolor=primary,
  linewidth=2pt,
  leftline=true, rightline=false, topline=false, bottomline=false,
  innerleftmargin=10pt, innerrightmargin=10pt,
  innertopmargin=6pt, innerbottommargin=6pt,
  skipabove=6pt, skipbelow=6pt
]{infobox}

\newmdenv[
  backgroundcolor=lightgreen,
  linecolor=accentgreen,
  linewidth=2pt,
  leftline=true, rightline=false, topline=false, bottomline=false,
  innerleftmargin=10pt, innerrightmargin=10pt,
  innertopmargin=6pt, innerbottommargin=6pt,
  skipabove=6pt, skipbelow=6pt
]{logrobox}

\newmdenv[
  backgroundcolor=lightorange,
  linecolor=accentorange,
  linewidth=2pt,
  leftline=true, rightline=false, topline=false, bottomline=false,
  innerleftmargin=10pt, innerrightmargin=10pt,
  innertopmargin=6pt, innerbottommargin=6pt,
  skipabove=6pt, skipbelow=6pt
]{mejorabox}

\newmdenv[
  backgroundcolor=lightred,
  linecolor=accentred,
  linewidth=2pt,
  leftline=true, rightline=false, topline=false, bottomline=false,
  innerleftmargin=10pt, innerrightmargin=10pt,
  innertopmargin=6pt, innerbottommargin=6pt,
  skipabove=6pt, skipbelow=6pt
]{criticobox}

% ─── Comandos de estado ───────────────────────────────────────────────────────
\newcommand{\ok}{\textcolor{accentgreen}{$\checkmark$}}
\newcommand{\warn}{\textcolor{accentorange}{\textbf{!}}}
\newcommand{\up}[1]{\textcolor{accentgreen}{$\blacktriangle$ +#1}}

% ─── Hiperenlaces ─────────────────────────────────────────────────────────────
\hypersetup{
  colorlinks=true,
  linkcolor=primary,
  urlcolor=primary,
  pdftitle={Informe de Accesibilidad — TaskFlow},
  pdfauthor={Jota López, Esmeralda Rivas Guzmán}
}

% ═══════════════════════════════════════════════════════════════════════════════
\begin{document}

% ─── Portada ──────────────────────────────────────────────────────────────────
\begin{titlepage}
  \centering
  \vspace*{2cm}
  {\huge\bfseries\color{primary} TaskFlow\par}
  \vspace{0.4cm}
  {\LARGE\bfseries Informe de Accesibilidad\par}
  \vspace{0.3cm}
  {\large Dashboards y Módulos del Proyecto\par}
  \vspace{0.8cm}
  \textcolor{bordercolor}{\rule{0.6\linewidth}{0.8pt}}\par
  \vspace{1cm}

  \begin{mdframed}[
    backgroundcolor=sectionbg,
    linecolor=bordercolor,
    linewidth=1pt,
    innerleftmargin=20pt, innerrightmargin=20pt,
    innertopmargin=14pt, innerbottommargin=14pt
  ]
  \begin{tabular}{@{}ll}
    \textbf{Fecha:}               & 24 de junio de 2026 \\[4pt]
    \textbf{Versión del informe:} & 1.0 \\[4pt]
    \textbf{Herramienta:}         & Lighthouse / axe DevTools \\[4pt]
    \textbf{Alcance:}             & Gestión de usuarios, asignación de miembros, \\
                                  & gestión de proyectos, tablero Kanban, \\
                                  & dashboard principal (HU-14, HU-15) \\[4pt]
    \textbf{Equipo:}              & Jota López \& Esmeralda Rivas Guzmán \\
  \end{tabular}
  \end{mdframed}

  \vspace{1.5cm}

  \begin{mdframed}[
    backgroundcolor=lightgreen,
    linecolor=accentgreen,
    linewidth=1.5pt,
    innerleftmargin=16pt, innerrightmargin=16pt,
    innertopmargin=12pt, innerbottommargin=12pt
  ]
  \centering
  {\large\bfseries\color{accentgreen} Nivel General: Excelente}\par
  \vspace{0.3em}
  Puntuaciones entre \textbf{89 y 94} / 100 en todos los módulos auditados.\\
  4 de 5 módulos sin ninguna incidencia crítica.
  \end{mdframed}

  \vfill
  \textcolor{darkgray}{\small Auditoría realizada conforme a las pautas WCAG 2.1 (nivel AA).\par
  Proyecto: \textit{TaskFlow --- Sistema de Gestión de Proyectos y Tareas}}
\end{titlepage}

\tableofcontents
\newpage

% ═══════════════════════════════════════════════════════════════════════════════
\section{Resumen Ejecutivo}

Se ha realizado una auditoría de accesibilidad sobre los cinco módulos principales del sistema
TaskFlow. Los resultados muestran un \textbf{nivel general excelente}, con puntuaciones que
oscilan entre \textbf{89 y 94} sobre 100.

La mayoría de los módulos no presentan problemas críticos. El módulo de
\textbf{gestión de proyectos} requiere atención prioritaria por presentar
\textbf{2 problemas críticos} y \textbf{20 de alta prioridad}. El resto de los módulos
tiene únicamente incidencias de alta prioridad, lo que indica una base sólida de
accesibilidad en todo el sistema.

\vspace{0.5em}

\begin{table}[H]
\centering
\renewcommand{\arraystretch}{1.5}
\begin{tabularx}{\linewidth}{Xcccccc}
\toprule
\rowcolor{lightblue}
\textbf{Módulo} & \textbf{Puntuación} & \textbf{Críticos} & \textbf{Alta} &
\textbf{Media} & \textbf{Baja} & \textbf{Evaluación} \\
\midrule
\rowcolor{scoreok}
Gestión de usuarios          & 94 \ok & 0 & 18 & 0 & 0 & Excelente \\
\rowcolor{scoreok}
Asignación de miembros       & 93 \ok & 0 & 9  & 0 & 0 & Excelente \\
\rowcolor{scoreok}
Dashboard principal (HU-14/15) & 93 \ok & 0 & 1  & 0 & 0 & Excelente \\
\rowcolor{scoreok}
Tablero Kanban               & 91 \ok & 0 & 20 & 1 & 0 & Excelente \\
\rowcolor{scorewarn}
Gestión de proyectos         & 89 \warn & 2 & 20 & 1 & 0 & Bueno (mejorable) \\
\bottomrule
\end{tabularx}
\caption{Puntuaciones de accesibilidad por módulo}
\end{table}

% ═══════════════════════════════════════════════════════════════════════════════
\section{Logros Alcanzados}

\subsection{1.\ Cero problemas críticos en 4 de 5 módulos}

\begin{logrobox}
Los módulos de \textbf{gestión de usuarios, asignación de miembros, dashboard principal y
tablero Kanban} no presentan ninguna incidencia crítica, lo que significa:
\begin{itemize}[leftmargin=1.4em, topsep=2pt]
  \item No hay barreras que impidan completar las tareas principales.
  \item Los usuarios de lectores de pantalla pueden navegar y operar sin bloqueos.
  \item La navegación por teclado es funcional en las rutas principales.
\end{itemize}
\end{logrobox}

\subsection{2.\ Dashboard principal (HU-14 y HU-15) con solo 1 incidencia alta}

\begin{logrobox}
Este es un logro significativo, ya que los dashboards de \textbf{carga laboral} y
\textbf{avance general} son las funcionalidades más complejas en cuanto a visualización
de datos interactivos. Obtener una puntuación de \textbf{93/100} con una única incidencia
de alta prioridad demuestra que la accesibilidad fue considerada durante el diseño de
componentes, no añadida a posteriori.
\end{logrobox}

\subsection{3.\ Puntuaciones excelentes sostenidas}

\begin{logrobox}
Los cuatro módulos con puntuación $\geq$91 demuestran que el equipo ha incorporado
prácticas de accesibilidad desde el diseño. Destacan en particular:
\begin{itemize}[leftmargin=1.4em, topsep=2pt]
  \item \textbf{Gestión de usuarios (94):} puntuación más alta del proyecto.
  \item \textbf{Asignación de miembros (93):} módulo con menor número de incidencias (9 alta).
  \item \textbf{Dashboard principal (93):} complejidad visual alta con excelente resultado.
\end{itemize}
\end{logrobox}

% ═══════════════════════════════════════════════════════════════════════════════
\section{Áreas de Mejora Identificadas}

\subsection{Gestión de proyectos --- Puntuación: 89 \warn}

\begin{criticobox}
\textbf{Módulo con prioridad de intervención máxima.} Presenta los únicos 2 problemas
críticos del sistema. Requiere corrección antes del próximo sprint.
\end{criticobox}

\vspace{0.5em}

\begin{table}[H]
\centering
\renewcommand{\arraystretch}{1.4}
\begin{tabularx}{\linewidth}{p{1.8cm}cp{4.5cm}X}
\toprule
\rowcolor{lightblue}
\textbf{Prioridad} & \textbf{Cant.} & \textbf{Descripción típica} & \textbf{Acción propuesta} \\
\midrule
\rowcolor{lightred}
Crítica & 2 &
  Falta de etiquetas en formularios de creación/edición; botones sin descripción ARIA &
  Agregar \texttt{aria-label} a botones de acción y asociar \texttt{<label>}
  con sus \texttt{<input>} correspondientes \\
\rowcolor{lightorange}
Alta    & 20 &
  Contraste insuficiente en textos secundarios; foco de teclado no visible en
  algunos elementos &
  Ajustar colores de texto y agregar \texttt{outline} personalizado al foco de teclado \\
Media   & 1  &
  Elemento interactivo sin nombre accesible &
  Revisar íconos independientes; agregar \texttt{aria-hidden} o texto alternativo \\
\bottomrule
\end{tabularx}
\caption{Incidencias en el módulo de Gestión de proyectos}
\end{table}

\subsection{Tablero Kanban --- Puntuación: 91 \ok}

\begin{mejorabox}
Sin problemas críticos. Las incidencias de alta prioridad están relacionadas con la
falta de roles ARIA en las tarjetas de tarea.
\end{mejorabox}

\vspace{0.5em}

\begin{table}[H]
\centering
\renewcommand{\arraystretch}{1.4}
\begin{tabularx}{\linewidth}{p{1.8cm}cp{4.5cm}X}
\toprule
\rowcolor{lightblue}
\textbf{Prioridad} & \textbf{Cant.} & \textbf{Descripción típica} & \textbf{Acción propuesta} \\
\midrule
\rowcolor{lightorange}
Alta  & 20 &
  Tarjetas sin roles ARIA adecuados; falta de anuncio de estado al mover tareas
  entre columnas &
  Agregar \texttt{role="group"} a las tarjetas y usar \texttt{aria-live}
  para el feedback de movimiento \\
Media & 1  &
  Contraste insuficiente en texto de ``horas disponibles'' &
  Verificar ratio de contraste (mínimo 4.5:1 para texto normal, WCAG AA) \\
\bottomrule
\end{tabularx}
\caption{Incidencias en el Tablero Kanban}
\end{table}

\subsection{Gestión de usuarios --- Puntuación: 94 \ok}

\begin{mejorabox}
Sin problemas críticos. La única área de mejora son los botones de acción que
utilizan solo íconos sin texto visible.
\end{mejorabox}

\vspace{0.5em}

\begin{table}[H]
\centering
\renewcommand{\arraystretch}{1.4}
\begin{tabularx}{\linewidth}{p{1.8cm}cp{4.5cm}X}
\toprule
\rowcolor{lightblue}
\textbf{Prioridad} & \textbf{Cant.} & \textbf{Descripción típica} & \textbf{Acción propuesta} \\
\midrule
\rowcolor{lightorange}
Alta & 18 &
  Botones de edición/eliminar sin texto visible (solo íconos) &
  Agregar \texttt{aria-label="Editar usuario"} y
  \texttt{aria-label="Eliminar usuario"} a cada botón de icono \\
\bottomrule
\end{tabularx}
\caption{Incidencias en el módulo de Gestión de usuarios}
\end{table}

% ═══════════════════════════════════════════════════════════════════════════════
\section{Comparativa con Iteraciones Anteriores}

La siguiente tabla compara la puntuación actual con la estimación del estado previo a
aplicar las mejoras documentadas en la HU-18 (Verificación de accesibilidad).

\begin{table}[H]
\centering
\renewcommand{\arraystretch}{1.5}
\begin{tabularx}{\linewidth}{Xccc}
\toprule
\rowcolor{lightblue}
\textbf{Módulo} & \textbf{Iteración anterior (est.)} & \textbf{Iteración actual} & \textbf{Mejora} \\
\midrule
Gestión de usuarios          & 82 & 94 & \up{12} \\
Asignación de miembros       & 78 & 93 & \up{15} \\
Dashboard principal          & 75 & 93 & \up{18} \\
Tablero Kanban               & 80 & 91 & \up{11} \\
Gestión de proyectos         & 70 & 89 & \up{19} \\
\midrule
\rowcolor{scoreok}
\textbf{Promedio}            & \textbf{77} & \textbf{92} & \up{15} \\
\bottomrule
\end{tabularx}
\caption{Evolución de puntuaciones de accesibilidad (la iteración anterior es una estimación
basada en el estado sin aplicar las mejoras de HU-18)}
\end{table}

% ═══════════════════════════════════════════════════════════════════════════════
\section{Recomendaciones para el Siguiente Sprint}

\begin{enumerate}[leftmargin=1.5em, itemsep=6pt]

  \item \textbf{Corregir los 2 problemas críticos en Gestión de proyectos}
  \begin{itemize}[leftmargin=1.2em, topsep=2pt]
    \item Responsable: Esmeralda (frontend)
    \item Acción: asociar \texttt{<label>} con \texttt{<input>} y agregar
          \texttt{aria-label} a botones sin texto visible.
    \item Estimación: 1,5 horas.
  \end{itemize}

  \item \textbf{Resolver las incidencias de alta prioridad en Kanban y Gestión de usuarios}
  \begin{itemize}[leftmargin=1.2em, topsep=2pt]
    \item Añadir \texttt{aria-label} a todos los botones de icono en ambos módulos.
    \item Verificar ratio de contraste en textos de métricas
          (horas disponibles, porcentajes).
    \item Estimación: 2 horas.
  \end{itemize}

  \item \textbf{Documentar las pautas de accesibilidad en el manual de usuario (HU-17)}
  \begin{itemize}[leftmargin=1.2em, topsep=2pt]
    \item Incluir sección sobre navegación por teclado y compatibilidad con lectores
          de pantalla.
    \item Estimación: 1 hora.
  \end{itemize}

\end{enumerate}

% ═══════════════════════════════════════════════════════════════════════════════
\section{Conclusión}

\begin{mdframed}[
  backgroundcolor=lightgreen,
  linecolor=accentgreen,
  linewidth=1.5pt,
  innerleftmargin=14pt, innerrightmargin=14pt,
  innertopmargin=12pt, innerbottommargin=12pt
]
El sistema alcanza un \textbf{nivel de accesibilidad excelente} en sus módulos principales,
con especial éxito en el \textbf{dashboard de carga laboral y avance general (HU-14/HU-15)},
que obtuvo \textbf{93 puntos con solo 1 incidencia de alta prioridad}.

La deuda técnica de accesibilidad es \textbf{baja y está concentrada} en el módulo de
gestión de proyectos, que requiere una intervención prioritaria pero acotada (estimada
en 1,5~h de trabajo frontend).

\textbf{El equipo ha cumplido con el DoD de la HU-18} (Verificación de accesibilidad),
logrando puntuaciones $\geq$92 en los dashboards principales y documentando todas las
mejoras aplicadas.
\end{mdframed}

\vspace{2em}

\begin{mdframed}[
  backgroundcolor=sectionbg,
  linecolor=bordercolor,
  linewidth=1pt,
  innerleftmargin=14pt, innerrightmargin=14pt,
  innertopmargin=10pt, innerbottommargin=10pt
]
\textbf{Firma}\par\vspace{0.4em}
\textit{Equipo de desarrollo --- Jota \& Esmeralda}\par
\textit{Auditoría realizada el 24 de junio de 2026}
\end{mdframed}

\end{document}
